% !PRINTED 205
\weberpage{205}
%<<CONT
\[
\psi_1 = a_1'\xi^2 - a_3'\eta^2
\]
\[
= a_0\,\frac{(\beta-\gamma)(\beta-\delta)(x-\alpha)^2 - (\alpha-\gamma)(\alpha-\delta)(x-\beta)^2}{\alpha-\beta}.
\]
Darin aber lässt sich $\alpha - \beta$ im Zähler und Nenner wegheben, und man kann
$\psi_1$ in zwei Formen darstellen; um $\psi_2$ und $\psi_3$ zu erhalten, braucht man
dann nur $\beta$, $\gamma$, $\delta$ cyklisch zu vertauschen.

Man findet so
\[
\begin{aligned}
\frac{\psi_1}{a_0} &= (\gamma-\beta)(x-\alpha)(x-\delta) - (\alpha-\delta)(x-\beta)(x-\gamma) \\
&= (\delta-\beta)(x-\alpha)(x-\gamma) - (\alpha-\gamma)(x-\beta)(x-\delta)
\end{aligned}
\]
\weq{6}{
\begin{aligned}
\frac{\psi_2}{a_0} &= (\delta-\gamma)(x-\alpha)(x-\beta) - (\alpha-\beta)(x-\gamma)(x-\delta) \\
&= (\beta-\gamma)(x-\alpha)(x-\delta) - (\alpha-\delta)(x-\gamma)(x-\beta) \\
\frac{\psi_3}{a_0} &= (\beta-\delta)(x-\alpha)(x-\gamma) - (\alpha-\gamma)(x-\delta)(x-\beta) \\
&= (\gamma-\delta)(x-\alpha)(x-\beta) - (\alpha-\beta)(x-\delta)(x-\gamma).
\end{aligned}
}
Die Grössen $\psi_1^2$, $\psi_2^2$, $\psi_3^2$ kann man auffassen als die Wurzeln
einer cubischen Gleichung, deren Co\"efficienten Covarianten sind. Man erhält
diese Gleichung leicht aus (4) in der Form
\[
\left(z + \frac{H+4a_2'f}{3}\right)\left(z + \frac{H+4a_2''f}{3}\right)\left(z + \frac{H+4a_2'''f}{3}\right) = 0.
\]
Diese Gleichung ergiebt nach \S.\,65, (12) und mit Benutzung des Ausdruckes (5)
für $T^2$
\weq{7}{
z^3 + Hz^2 + \tfrac{1}{3}(H^2 - 16Af^2)z - T^2 = 0.
}
Es ist noch von Interesse, die Discriminante dieser Gleichung zu bilden. Man
erhält sie am einfachsten aus dem Ausdruck
\[
\varDelta = (\psi_1^2 - \psi_2^2)^2(\psi_1^2 - \psi_3^2)^2(\psi_2^2 - \psi_3^2)^2,
\]
wenn man die Ausdrücke (4) einsetzt und dann für $(a_2'-a_2'')^2(a_2'-a_2''')^2(a_2''-a_2''')^2$
die Discriminante der Gleichung \S.\,65, (12)
\[
27(4A^3 - B^2) = 3^6 D
\]
setzt, wo dann $D$ die Discriminante von $f$ ist. Man erhält so
\weq{8}{
\varDelta = 2^{12} D f^6.
}
